% Proof-ready contributions by /root/programme_audit.
% Run JAC1964-20261121; GPT-6 Astra Pro; effort setting not exposed.
% All statements are rational, unless explicitly stated otherwise.

\section{Incidence Picard varieties and the next theorem of the programme}

For a smooth projective variety $V$ of dimension $d$ over an algebraically
closed field, write $A^p(V)=\mathrm{CH}^p(V)_{\mathrm{alg}}$ and
\[
 I^p(V)=\{z\in A^p(V):\Gamma_*z=0\text{ in }\mathrm{Pic}^0(T)
 \text{ for every smooth projective }T
 \text{ and }\Gamma\in\mathrm{CH}^{d-p+1}(V\times T)\}.
\]
Saito's higher Picard variety, denoted $P^p(V)$ here, represents
$A^p(V)/I^p(V)$ on algebraically closed points
\cite[Theorem~(3.6)(i)]{Saito1979}.
This is the incidence construction suggested by cote~44, scans~2 and~9.
It must not be silently substituted for the universal regular algebraic
representative, a different universal problem.  Saito supplies duality
\emph{up to isogeny}, not an assertion of a canonical integral principal
polarization \cite[Proposition~(4.4)]{Saito1979}.
Rational families below mean rational homomorphisms of abelian varieties,
obtained by clearing the denominators of the cycles first.  They do not
mean that one tensors the abstract group of geometric points with
$\mathbf Q$; that operation would lose the construction over
$\overline{\mathbf F}_p$, where all abelian-variety points are torsion.

\begin{lemma}[JAC-P00: rational descent of the incidence construction]
\label{jac:P00}
Let $F$ be any field and let $V/F$ be smooth and projective.  The
incidence Picard variety of $V_{\overline F}$, together with its quotient
map on algebraically trivial cycles, has a model over $F$ in the isogeny
category.  If $T/F$ is abelian and $u:T\to P^p(V)$ is a rational
homomorphism over $F$, there is a cycle
$Z\in\mathrm{CH}^p(T\times V)_{\mathbf Q}$ whose family map is $u$.
\end{lemma}

\begin{proof}
Saito's quotient is canonical: its universal property among Picard
homomorphisms gives a unique isomorphism from each conjugate quotient
to the quotient of $V_{\overline F}$.  These isomorphisms satisfy the
cocycle identity.  The variety, its group law, the finite correspondence
data witnessing that its quotient map is a Picard homomorphism, and
one surjective generating family are defined over a finite extension
of $F$.  The latter can be taken from Saito's Poincar\'e cycle, whose
family map is a positive integer times the identity.  An automorphism
fixing these data fixes the quotient map on that generating family.
Uniqueness of the comparison isomorphism, and surjectivity of the
family map, then force that comparison to be the identity.  This
proves, using finite algebraic data rather than the set of all cycles,
that the canonical descent datum is trivial on an open subgroup of
the absolute Galois group.

Here are the details that also cover an imperfect $F$.  Enlarge a
finite separable extension to a finite Galois extension $L/F$, and
choose $r\ge0$ so that the preceding finite field of definition lies
in $L^{1/p^r}$ in characteristic $p$; take $r=0$ in characteristic
zero.  Raising the coefficients of the variety and its group law to
$p^r$ gives an abelian variety $Q/L$, and relative Frobenius gives an
isogeny $f:P_{\overline F}\to Q_{\overline F}$.
Transport the canonical descent datum by $f$ and its rational inverse.
For $\sigma\in\operatorname{Gal}(L/F)$ this gives
\[
 c_\sigma:{}^\sigma Q\longrightarrow Q
 \quad\text{in the rational category},\qquad
 c_{\sigma\tau}=c_\sigma\,{}^\sigma c_\tau.
\]
These maps are defined over $L$: homomorphisms of abelian varieties
acquire no new elements under purely inseparable extensions, and the
remaining Galois invariance follows from the cocycle identity and the
triviality of the datum on automorphisms fixing $L$.  For the former
standard rigidity fact, see \cite[Example~4.2.14, p.~38]{ConradAV}.

Rational Galois descent is effective by the following explicit
construction.  Let $G=\operatorname{Gal}(L/F)$ and
$R=\operatorname{Res}_{L/F}Q$.  Over $L$ identify
$R_L=\prod_{\sigma\in G}{}^\sigma Q$.  The maps
\[
 i:Q\longrightarrow R_L,\quad
 i(x)=(c_\sigma^{-1}x)_\sigma,
 \qquad
 q:R_L\longrightarrow Q,\quad
 q((x_\sigma))=\frac1{|G|}\sum_\sigma c_\sigma x_\sigma
\]
satisfy $qi=1$.  The cocycle identity makes $e=iq$ Galois invariant,
so $e\in\operatorname{End}^0_F(R)$.  For a positive integer $N$ with
$Ne$ integral, its image is an abelian variety over $F$ whose base
change is isomorphic to $P$ in the isogeny category.  This supplies the
required rational model and its transported quotient map.  Uniqueness
in the rational category follows from the same quotient universal
property.

Finally clear a positive denominator $d$ so that $du$ is a morphism,
pull a normalized rational Poincar\'e cycle back by $du$ over
$\overline F$, and divide the resulting cycle by $d$.
The resulting family is defined over some finite
extension $F'/F$.  Put $d_s=[F':F]_{\rm sep}$ and
$d_i=[F':F]_{\rm insep}$, and let
$\pi:(T\times V)_{F'}\to T\times V$.
Define $Z=(d_sd_i)^{-1}\pi_*Z_{F'}$.  After base change to
$\overline F$, flat base change for Chow groups expresses this as
the average of the $d_s$ conjugate cycles, each with multiplicity
$d_i$.  The multiplicity is the length of the corresponding local
factor of $F'\otimes_F\overline F$.  Every conjugate induces $u$,
because $u$ is defined over $F$ for the compatible quotient descent
datum.  Hence $Z$ induces $u$.  This explains explicitly both the
separable conjugates and the inseparable factor; it does not treat a
purely inseparable extension as if it had extra automorphisms.
\end{proof}

\begin{lemma}[JAC-P01: detection by the odd realization]
\label{jac:P01}
Let $V,W$ be smooth projective varieties over an algebraically closed
field, and let
$\Gamma\in\mathrm{CH}^{\dim V-p+q}(V\times W)_{\mathbf Q}$.
Fix a prime $\ell$ invertible in the field.  If the action of $\Gamma$
from $H^{2p-1}(V,\mathbf Q_\ell)$ to
$H^{2q-1}(W,\mathbf Q_\ell(q-p))$ is zero, then the induced rational
homomorphism $P^p(V)\longrightarrow P^q(W)$ is zero.

More generally, for an abelian variety $T$ and
$Z\in\mathrm{CH}^q(T\times W)_{\mathbf Q}$, the homomorphism
$T\longrightarrow P^q(W)$ defined by $t\mapsto Z_t-Z_0$ depends only
on the K\"unneth component
\[
 [Z]_{1,2q-1}\in H^1(T,\mathbf Q_\ell)\otimes
 H^{2q-1}(W,\mathbf Q_\ell(q)).
\]
In particular, it is zero if that component is zero.
\end{lemma}

\begin{proof}
For the second assertion, test $Z_t-Z_0$ by an arbitrary incidence
correspondence
$\Xi\in\mathrm{CH}^{\dim W-q+1}(W\times S)$.
The composition $\Xi\circ Z$ is a divisor class on $T\times S$.
Its component in $H^1(T)\otimes H^1(S)$ depends only on
$[Z]_{1,2q-1}$, by the K\"unneth formula and integration over $W$.
A divisor correspondence on $T\times S$ gives a homomorphism
$T\to\mathrm{Pic}^0(S)$; its rational Tate-module action is precisely
the mixed degree-one component.  Faithfulness of the rational Tate
module on homomorphisms of abelian varieties therefore says that this
homomorphism is zero when the indicated mixed component is zero.
Consequently every $Z_t-Z_0$ is annihilated by every incidence tester.
It is zero in $A^q(W)/I^q(W)$ by definition, proving the assertion.
The argument applies to rational cycles after clearing denominators.

Every algebraically trivial cycle is a finite sum of differences of
fibres of families parametrized by smooth projective curves.
For one such family
$U\in\mathrm{CH}^{p}(C\times V)$, its only relevant odd component
gives the map $H^1(C)\to H^{2p-1}(V)(p-1)$.
Composition with $\Gamma$ is zero by the hypothesis.  Composing further
with each incidence tester on $W$ gives a divisor correspondence on
$C\times S$ with zero mixed component.  The induced
$\operatorname{Alb}(C)\to\mathrm{Pic}^0(S)$ is zero, by the same
faithfulness argument.  Thus $\Gamma_*$ sends every element of
$A^p(V)$ into $I^q(W)$, and its action on $P^p(V)$ is zero.
\end{proof}

The lemma does not assume the standard conjecture~D.  The use of
cohomology concerns correspondences acting on families of cycles;
it does not identify homological and incidence equivalence of individual
cycles.

\begin{theorem}[JAC-P02: Theorem~6 for incidence Picard varieties]
\label{jac:P02}
Let $X$ be smooth projective of dimension $n$ over an algebraically
closed field, and assume that its normalized Lefschetz lowering
operator is algebraic.  For a smooth hyperplane section $j:Y\hookrightarrow X$,
restriction induces
\[
 j^*:P^p(X)\longrightarrow P^p(Y).
\]
It is an isogeny if $2p\le n-1$, has finite kernel if $2p=n$, and
is surjective if $2p\ge n+1$.  Each assertion means the corresponding
split statement in the rational category of abelian varieties, hence
the asserted statement for the original integral homomorphism.

Suppose further that $Y=Y_\eta$ is the generic member of a pencil with
smooth codimension-two base locus, and put $K=k(t)$.  In the critical
case $n=2p$, the image of $P^p(X)_K$ is the image of the $K/k$-trace
of $P^p(Y_\eta)$, up to isogeny.  Thus the rational form of all clauses
of cote~44, scan~4, Theorem~6, holds whenever the Lefschetz standard
conjecture holds for $X$.  In particular, it holds for abelian $X$
in every characteristic.
\end{theorem}

\begin{proof}
Write $L$ for multiplication by the hyperplane class and $\Lambda$
for the normalized lowering operator: it sends $L^a v$ to
$L^{a-1}v$ for a primitive vector $v$ and $a>0$, and kills $v$.
The standard conjecture supplies an algebraic correspondence for this
operator.  The correspondence
\[
 R=\Lambda_X\circ j_*
\]
has degree zero and induces $P^p(Y)\to P^p(X)$.
Put $r=2p-1$.
For $r<n$, no Lefschetz string terminates in degree $r$, so
\[
 Rj^*=\Lambda_X L_X=\operatorname{id}\quad\text{on }H^r(X).
\]
By Lemma~\ref{jac:P01}, the same identity holds on $P^p(X)$.
If $r<n-1$, weak Lefschetz makes $j^*:H^r(X)\to H^r(Y)$ an
isomorphism, hence $j^*R$ is also the identity on $H^r(Y)$ and on
$P^p(Y)$.  This proves the first assertion and, when $r=n-1$,
the assertion of finite kernel.

For $r\ge n$, the Gysin map
$j_*:H^r(Y)\to H^{r+2}(X)(1)$ is an isomorphism by the dual of
weak Lefschetz.  No Lefschetz string begins in degree $r+2>n$,
so $L_X\Lambda_X=\operatorname{id}$ there.  For $y\in H^r(Y)$,
\[
 j_*j^*Ry=L_X\Lambda_Xj_*y=j_*y.
\]
Injectivity of $j_*$ gives $j^*R=\operatorname{id}$ on $H^r(Y)$;
Lemma~\ref{jac:P01} transfers this identity to $P^p(Y)$.
These are precisely the algebraic inverse mechanisms of cote~16,
scan~33, formula~B(4), and scans~40 and~45, clauses~(iv) and~(iv bis).

Here is a cycle-level proof of the fixed-part assertion, so no inference
from invariant cohomology to algebraicity of a correspondence is left
unstated.  Write $p=m$, $n=2m$, let $B$ be the base locus, and let
\[
 b:\widetilde X=\operatorname{Bl}_B X\longrightarrow X,
 \qquad e:E\hookrightarrow\widetilde X,
 \qquad q:E\longrightarrow B
\]
be the blow-up, exceptional divisor, and its projection.  Let $T/k$
be an abelian variety and
$u:T_K\longrightarrow P^m(Y_\eta)$ a homomorphism.  Saito's
Poincar\'e $m$-cycle supplies, after division by its positive integer
multiplier, a rational cycle on $P^m(Y_\eta)\times Y_\eta$
inducing the identity \cite[\S4.1]{Saito1979}.  Lemma~\ref{jac:P00}
descends its pullback by $u$ to a rational cycle
\[
 Z_\eta\in\mathrm{CH}^m(T_K\times Y_\eta)_{\mathbf Q}
\]
inducing $u$.  In particular the construction is valid even when
$K=k(t)$ is imperfect.  Subtracting $Z_0$ normalizes
the family at the identity of $T$.

Close a cycle representative in $T\times\widetilde X$.
The codimension-two blow-up formula gives
\[
 \overline Z=(1_T\times b)^*Z_X+
 (1_T\times e)_*(1_T\times q)^*Z_B,
\]
where
$Z_X\in\mathrm{CH}^m(T\times X)_{\mathbf Q}$ and
$Z_B\in\mathrm{CH}^{m-1}(T\times B)_{\mathbf Q}$ are defined over
$k$.  Write $k_B:B\hookrightarrow X$ and $i:B_K\hookrightarrow Y_\eta$.
The mixed component relevant to $Z_B$ is
$H^1(T)\otimes H^{n-3}(B)$.
Weak Lefschetz identifies $H^{n-3}(X)$ with $H^{n-3}(B)$.
An algebraic inverse is
\[
 S=\Lambda_X^2\circ(k_B)_*.
\]
Indeed, $(k_B)_*k_B^*=L_X^2$ and $\Lambda_X^2L_X^2$ is the
identity on degree $n-3$.  Define an algebraic rational cycle over $k$
by
\[
 W_X=Z_X+L_X\circ S\circ Z_B\in
 \mathrm{CH}^m(T\times X)_{\mathbf Q}.
\]
For $x\in H^{n-3}(X)$ the projection formula gives
\[
 i_*k_B^*x=i_*i^*j_\eta^*x=L_{Y_\eta}j_\eta^*x
 =j_\eta^*L_Xx.
\]
It follows that $Z_\eta$ and $j_\eta^*W_X$ have the same
$H^1(T)\otimes H^{n-1}(Y_\eta)$ component.  Lemma~\ref{jac:P01}
says they induce the same rational homomorphism from $T_K$ to
$P^m(Y_\eta)$.  The latter factors through $P^m(X)_K$ since $W_X$
is defined over $k$.  When $m=1$ the exceptional term has no mixed
component, since $H^{-1}(B)=0$, and the same conclusion follows.

Every homomorphism from a constant abelian variety has therefore image
in $\operatorname{Im}j_\eta^*$.  Conversely, this image is itself the
image of a homomorphism from the constant variety $P^m(X)_K$.
The universal property of the $K/k$-trace proves the asserted equality
of images up to isogeny, using the rational $K$-model of
Lemma~\ref{jac:P00}.
\end{proof}

\section{An explicit test of the primitive sign}

\begin{proposition}[JAC-P03: the uncorrected correspondence is indefinite]
\label{jac:P03}
There exist normalized algebraic families of codimension-two cycles
on an abelian threefold for which the divisor correspondence of
cote~44, scan~9, is nondegenerate and indefinite.  Neither quotienting
by its kernel nor multiplication by a common sign makes it a
polarization.  The primitive qualification in scan~8 is essential.
\end{proposition}

\begin{proof}
Let $E$ be an elliptic curve, set $X=E_1\times E_2\times E_3$,
write $a_i$ for the pullback of $[0]\in\mathrm{CH}^1(E_i)$, and use
$h=a_1+a_2+a_3$ as polarization.  On $E\times E$ put
\[
 P=[\Delta_E]-[E\times0]-[0\times E].
\]
As a correspondence $P$ is the symmetric degree-one Chow--K\"unneth
projector, hence ${}^tP\circ P=P$.
As a divisor correspondence it induces
\[
 E\longrightarrow\widehat E,
 \qquad t\longmapsto\mathcal O_E(t-0)=-\lambda_{\mathcal O_E(0)}(t),
\]
where the polarization convention is
$\lambda_L(t)=t_t^*L\otimes L^{-1}$.
Thus the sign calibration is fixed, rather than hidden in a convention.

Define two families on $E_t\times X$ by
\[
 z_\pm=P_{t,1}\,(a_2\pm a_3).
\]
Their incidence divisor correspondences are
\[
 D_\pm=\deg_{E_2\times E_3}(a_2\pm a_3)^2\,P
       =\pm2P,
\]
and the mixed incidence correspondence between the two families
is zero because $(a_2+a_3)(a_2-a_3)=0$.
Use independent parameters for the two families, obtaining a family
over $T=E\times E$.  Its symmetric homomorphism $T\to\widehat T$
has diagonal blocks $-2\lambda_{\mathcal O(0)}$ and
$+2\lambda_{\mathcal O(0)}$.  It is an isogeny, but its associated
rational line bundle restricts to a negative-degree line bundle on one
elliptic factor and a positive-degree one on the other.  It cannot be
ample.  Its negative cannot be ample either.  Passing through a finite
quotient does not change this conclusion, since the pullback of an
ample line bundle by an isogeny is ample.

The degree-one parameter components make the primitive issue explicit.
For $\alpha\in H^1(E_1)$, they have images
\[
 u_\pm(\alpha)=\alpha(a_2\pm a_3).
\]
One has $u_+(\alpha)=L_X\alpha$ and $L_Xu_-(\alpha)=0$.
Moreover, for $\alpha,\beta\in H^1(E_1)$,
\[
 \int_Xu_+(\alpha)u_+(\beta)=2\int_E\alpha\beta,
 \qquad
 \int_Xu_-(\alpha)u_-(\beta)=-2\int_E\alpha\beta,
\]
and the cross pairing vanishes.  The lower Lefschetz and primitive
pieces require different signs.  In the divisor convention above,
the primitive $m=2$ correspondence $D_-=-2P$ already gives the
positive homomorphism $2\lambda_{\mathcal O(0)}$.
\end{proof}

\section{Primitive middle cycles need not reach the incidence Jacobian}

\begin{proposition}[JAC-P04: an abelian-fourfold domain test for H9]
\label{jac:P04}
Let $X=E_1\times E_2\times E_3\times E_4$ over $\mathbf C$,
$a_i=[0]$ pulled back from $E_i$, and $h=a_1+a_2+a_3+a_4$.
For a sufficiently high-degree Lefschetz pencil relative to $h$, let
$Y_\eta$ be its geometric generic member.  The codimension-two cycle
\[
 \alpha=(a_1-a_2)(a_3-a_4)
\]
satisfies $h\alpha=0$ and $\deg\alpha^2=4$.  Its restriction to
$Y_\eta$ is homologically trivial, but no nonzero multiple of that
restriction is algebraically equivalent to zero.
\end{proposition}

\begin{proof}
The relations $a_i^2=0$ give
\[
 h\alpha=0,
 \qquad
 \alpha^2=(-2a_1a_2)(-2a_3a_4)=4a_1a_2a_3a_4.
\]
In particular $\alpha$ has nonzero cohomology class.  Since the Gysin
map $H^4(Y_\eta)\to H^6(X)$ is an isomorphism by weak Lefschetz
and duality, the equality $j_*j^*\alpha=d\,h\alpha=0$ gives
$j^*\alpha=0$ in cohomology, for the degree-$d$ pencil.

Apply SGA~7~II, Expos\'e~XX, Theorem~4.1
\cite[printed p.~355]{SGA7II}.  Its hard-Lefschetz hypothesis holds
for $X$.  Its algebraic-retraction hypothesis~B is supplied by the
algebraic Lefschetz operator, as noted in Expos\'e~XX, \S3.4.
For sufficiently high degree the pencil satisfies condition~A by
Expos\'e~XVIII, Corollary~6.4.  In characteristic zero its variable
middle cohomology has level greater than one for sufficiently high
degree, by Expos\'e~XX, Remark~4.3 (printed p.~356).
The cited theorem says that a middle-dimensional cycle on $X$ with
algebraically trivial generic restriction has zero cohomology class
on $X$.  Applied to any nonzero multiple of $\alpha$, this contradicts
$\deg\alpha^2=4$.  Therefore no such multiple is algebraically
trivial on $Y_\eta$.
\end{proof}

This example identifies the enlargement between cote~44, scan~4,
Theorem~5, which explicitly requires algebraic triviality on $Y$,
and H9, which assumes only homological triviality.  The latter condition
does not provide a point of the incidence Picard variety.  It provides
a class in the Griffiths group in this example.  A general intermediate
Jacobian or normal function is a different object and does not remove
this distinction.

\begin{proposition}[JAC-P05: the customary reading of Theorem~7 is false]
\label{jac:P05}
If $\tau$-equivalence means that a nonzero integral multiple is
algebraically equivalent to zero, the unrestricted equality of numerical
and $\tau$-equivalence in cote~44, scan~6, fails already for complex
abelian threefolds of CM type, including one defined over a number field.
\end{proposition}

\begin{proof}
Take the Fermat quartic
\[
 C:\quad T_0^4+T_1^4=T_2^4
\]
and embed it in its Jacobian $J=J(C)$ using a rational base point.
Bloch's Theorem~(4.1), printed p.~104 of \cite{Bloch1984L}, states
that its Ceresa cycle $\Xi=[C]-[-1]_*[C]$ has infinite order modulo
algebraic equivalence.  This is the rational non-torsion assertion,
stronger than merely $\Xi\not\sim_{\rm alg}0$.
Bloch's Lemma~(3.1), printed p.~99, gives an isogeny
\[
 J\longrightarrow E\times E\times E',\qquad
 E:y^2=x^3-x,\quad E':y^2=x^3+x.
\]
Both elliptic curves have CM by $\mathbf Q(i)$, as is seen from
$(x,y)\mapsto(-x,iy)$; thus $J$ is of CM type.
On $H^4(J,\mathbf Q)$ the involution $[-1]^*$
acts as $(-1)^4=1$.  Thus this cycle is homologically trivial, hence
numerically trivial by compatibility of cycle classes and intersection.
Its infinite order modulo algebraic equivalence means that it is not
$\tau$-equivalent to zero.  This gives the required counterexample.
\end{proof}

Ceresa's original theorem also gives a very general genus-three
Jacobian example.  The explicit Bloch example has the advantage that
both the infinite-order statement and the CM isogeny have directly
verified primary-source locators.  A theorem making Hodge classes
algebraic says nothing that would force a homologically trivial
cycle such as $\Xi$ to become algebraically trivial.

For divisors the statement is correct in every characteristic:
SGA~6, Expos\'e~XIII, Theorem~4.6 identifies numerical triviality
with $\tau$-triviality; definitions are in \S\S4.1 and~4.4.
The supplied archive page does not itself define $\tau$.  Accordingly,
the preceding refutation is explicitly tied to this customary reading,
not to a guessed replacement of $\tau$ by homological equivalence.

\begin{lemma}[JAC-P06: the finite-field meaning of the last claim]
\label{jac:P06}
Over $k=\overline{\mathbf F}_p$, every algebraically trivial Chow class
is torsion.  Hence $\tau$-triviality means torsion in the Chow group,
and numerical equivalence equal to $\tau$-equivalence would identify
the rational Chow group with cycles modulo numerical equivalence.
For abelian varieties, and more generally geometric abelian-type
varieties, JAC-P15 proves that comparison using the accepted Tate
consequence of 032 and Kahn's additional established theorem.
\end{lemma}

\begin{proof}
An algebraically trivial cycle is a finite sum of images, under cycle
correspondences, of degree-zero zero-cycles on smooth projective curves.
For a curve $C/k$, its degree-zero Chow group is $J(C)(k)$.  Every
point and the curve are defined over a finite field, and each finite
group $J(C)(\mathbf F_q)$ is torsion.  Thus $J(C)(k)$ is torsion,
and so is each of its images under a correspondence.  A finite sum of
torsion elements is torsion.  If a nonzero multiple of a Chow class
is algebraically trivial, another nonzero multiple is therefore zero
in the Chow group.  The converse is immediate.
\end{proof}
